% !TeX spellcheck = en_US
%%DATE 2.4 12:01
%\chapter{Abstract-Deutsch}
%{

%abstrakt deutsch einfügen


%Die Diplomarbeit befasst sich mit der Entwicklung des Eishockey-Trainingssystems \glqq PuckVision\grqq, dessen Herzstück eine 0,80 m x 0,80 m große Kunststoffplatte ist, die als Spielfläche dient. Das Ziel ist es, den Puck so schnell wie möglich über aufleuchtende Spots zu bewegen. Jeder Spot besteht aus einer LED-Beleuchtung und einem Halleffektsensor, der präzise erkennt, ob der Puck erfolgreich über diesen geführt wird. Die LED visualisiert den Spot und zeigt die jeweilige Trainingsaufgabe an.
%Die Daten werden in Echtzeit über das Kommunikationsprotokoll I\textsuperscript{2}C an eine zentrale Steuereinheit, den Mikrocontroller \glqq ESP32 \grqq übertragen, die alle Sensoren, LEDs und Spielfunktionen synchronisiert und steuert. Spieler können ihre Fortschritte über eine begleitende Webseite verfolgen, die Daten wie den Punktestand oder eine Option zum Zurücksetzen zur Verfügung stellt. 